\chapter{Radial variation and exact monotonicity}\label{sec:varifold-exact-monotonicity}
Throughout this section $1\leq m<N$, $V$ is an $m$-varifold in an open $U\subset\mathbb R^N$, and $\mu=\|V\|$. Balls are open. The calculations develop \cite{Sim18}, Chapter 8, equations (3.3)--(3.6), pp.~240--241, including the endpoint convention needed for general varifolds.

\begin{definition}[Mass ratio and radial tilt]\label{def:radial-varifold-tilt}
For $\overline{B_R(a)}\subset U$, put
\[
 f_a(r)=r^{-m}\mu(B_r(a)),\qquad
 T_a(x,S)=\frac{|\pi_{S^\perp}(x-a)|^2}{|x-a|^{m+2}}\quad(x\ne a),
\]
and set $T_a(a,S)=0$. The normalized mass ratio is $f_a(r)/\omega_m$. Radial tilt vanishes when the displacement $x-a$ lies in the plane recorded by the varifold; this does not require that the varifold be rectifiable. Compare \cite{Sim18}, Chapter 8, (3.3)--(3.6), pp.~240--241.
\end{definition}

\begin{lemma}[Divergence of a radial test field]\label{lem:radial-varifold-divergence}
Write $r=|x-a|$. If $X(x)=\gamma(r)(x-a)$ is $C^1$, then for $r>0$
\[
 \operatorname{div}_S X=m\gamma(r)+r\gamma'(r)
 -\frac{\gamma'(r)}r|\pi_{S^\perp}(x-a)|^2.
\]
At $a$, if $\gamma$ is constant near zero, the divergence is $m\gamma(0)$.
\end{lemma}
\begin{proof}
For a unit vector $e\in S$, differentiation gives
$\langle D_eX,e\rangle=\gamma+\gamma'\langle x-a,e\rangle^2/r$.
Sum over an orthonormal basis of $S$ and use
$|\pi_S(x-a)|^2=r^2-|\pi_{S^\perp}(x-a)|^2$. The value at the center follows directly from the linear field there.
\end{proof}

\begin{proposition}[Exact stationary annular identity]\label{prop:stationary-annular-identity}
If $V$ is stationary and $0<\sigma<\rho<R$, then
\[
 f_a(\rho)-f_a(\sigma)
 =\int_{G_m(B_\rho(a)\setminus B_\sigma(a))}T_a(x,S)\,dV(x,S).
\]
The annulus includes its inner sphere and excludes its outer sphere. No assumption about the mass of either sphere is imposed.
\end{proposition}
\begin{proof}
First suppose both spheres have zero $\mu$ mass. Use the radial cutoff
\[
 \gamma(r)=\begin{cases}
 \sigma^{-m}-\rho^{-m},&0\leq r\leq\sigma,\\
 r^{-m}-\rho^{-m},&\sigma<r<\rho,\\
 0,&r\geq\rho.
 \end{cases}
\]
Smooth the two corners, keeping the derivative uniformly bounded and the support in a fixed compact subset of $U$. Away from the two spheres the fields and their derivatives converge. Dominated convergence in first variation is therefore valid. Lemma \ref{lem:radial-varifold-divergence} gives
\[
 0=m\sigma^{-m}\mu(B_\sigma)-m\rho^{-m}\mu(B_\rho)
   +m\int_{G_m(B_\rho\setminus B_\sigma)}T_a\,dV.
\]
For arbitrary radii choose continuity radii $\sigma_j\uparrow\sigma$ and $\rho_j\uparrow\rho$, with $\sigma_j<\rho_j$. Such choices exist because disjoint positive-mass spheres are at most countable in a compact ball. Open-ball masses converge from below. The annular indicators converge to that of $\{\sigma\leq|x-a|<\rho\}$; the integrands are bounded on a common annulus separated from $a$. Dominated convergence proves the identity, including both endpoint conventions.
\end{proof}

\begin{corollary}[Density defect and absence of atoms]\label{cor:stationary-density-defect}
For stationary $V$, $\mu(\{a\})=0$ and
\[
 f_a(\rho)-\omega_m\Theta^m(\mu,a)
 =\int_{G_m(B_\rho(a))}T_a\,dV<\infty\qquad(0<\rho<R).
\]
\end{corollary}
\begin{proof}
Theorem \ref{thm:stationary-monotonicity} bounds $\mu(B_r(a))$ by $Cr^m$ for small $r$, which excludes an atom since $m\geq1$. In Proposition \ref{prop:stationary-annular-identity} send $\sigma\downarrow0$. The left side converges by the definition of density, and monotone convergence applies on the right. The omitted fiber over $a$ has zero mass.
\end{proof}

\begin{corollary}[Equality and quantitative radial tilt]\label{cor:stationary-radial-equality}
If $0<\sigma<\rho<R$, then $f_a(\rho)=f_a(\sigma)$ precisely when $x-a\in S$ for $V$-almost every $(x,S)$ with $\sigma\leq|x-a|<\rho$. Moreover
\[
 \int_{G_m(B_\rho\setminus B_\sigma)}|\pi_{S^\perp}(x-a)|^2\,dV
 \leq\rho^{m+2}\bigl(f_a(\rho)-f_a(\sigma)\bigr).
\]
\end{corollary}
\begin{proof}
The integrand in Proposition \ref{prop:stationary-annular-identity} is nonnegative and its denominator is positive on the annulus. An integral of such a function vanishes exactly when the numerator vanishes almost everywhere. Multiplication by the upper bound $|x-a|^{m+2}\leq\rho^{m+2}$ gives the estimate. This controls radial tilt, not distance from a fixed plane.
\end{proof}

\begin{example}[A stationary varifold with zero density on its support]\label{ex:diffuse-stationary-varifold}
Fix an $m$-plane $P\subset\mathbb R^N$, $m<N$. The general varifold $V=\mathcal L^N\otimes\delta_P$ is stationary, has full support, and has zero $m$-density at every point.
\end{example}
\begin{proof}
In coordinates adapted to $P$, first variation is the integral over $\mathbb R^N$ of $\sum_{i=1}^m\partial_iX_i$, which vanishes by Fubini and compact support. Its normalized mass ratio is $(\omega_N/\omega_m)r^{N-m}$, tending to zero. Its weight is not $m$-rectifiable. Thus Corollary \ref{cor:integral-support-density} cannot be extended to arbitrary stationary varifolds.
\end{proof}

\begin{example}[A stationary varifold that charges a sphere]\label{ex:stationary-circle-rays}
In $\mathbb R^2$ put $e_\theta=(\cos\theta,\sin\theta)$ and define a general $1$-varifold by
\[
 V(\varphi)=\int_0^{2\pi}\varphi(e_\theta,T_{e_\theta}S^1)\,d\theta
 +\int_0^{2\pi}\int_1^\infty\varphi(r e_\theta,\mathbb R e_\theta)\,dr\,d\theta.
\]
It is stationary, $\|V\|(S^1)=2\pi$, and
\[
 \|V\|(B_r(0))=\begin{cases}0,&0<r\leq1,\\2\pi r,&r>1.\end{cases}
\]
Consequently the mass ratio of a general stationary varifold need not be continuous or absolutely continuous in radius.
\end{example}
\begin{proof}
Both terms define locally finite nonnegative measures on the Grassmann bundle. Integration by parts along the circle gives first variation $\int_0^{2\pi}X(e_\theta)\cdot e_\theta\,d\theta$. Along each outward ray it gives $-X(e_\theta)\cdot e_\theta$, so their sum vanishes. The ray endpoints have zero length measure. For $r>1$ their total mass in the ball is $2\pi(r-1)$, in addition to the circle mass. For $r\leq1$ the open ball contains neither term. Outside the circle the weight has planar density $1/|x|$, so this is not a rectifiable $1$-varifold. The circle supplies exactly the radial-tilt jump in Proposition \ref{prop:stationary-annular-identity}.
\end{proof}

\begin{proposition}[No stationary varifold supported in a compact interior set]\label{prop:no-compact-stationary-varifold}
If $V$ is stationary in $U$ and $\operatorname{spt}\|V\|$ is compactly contained in $U$, then $V=0$.
\end{proposition}
\begin{proof}
Choose a smooth compactly supported cutoff equal to one on a neighborhood of the support and multiply it by $x-a$ for any fixed $a$. Its tangential divergence is $m$ on the support. Stationarity gives $0=m\|V\|(U)$, so the weight and hence the varifold vanish. The compact-containment assumption rules out boundary terms at $\partial U$.
\end{proof}
